\chapter{Multi-factor diffusion models}
One-factor models are not perfect as they cannot generate all the shapes of empirically observed
yield curves. Specifically they cannot generate, they cannot generate trough and twists of the yield
curve and any changes in two yields or bond prices will be perfectly correlated (change in short
yield affects long yield).

In this presentation I will introduce the concept of Multi-factor models try to accommodate this
problem, by adding more state variables, which can affect the model both independent of, or
dependent on the short rate. We will consider the subclass of Affine Multi-factor models and
describe how one can price zero-coupon bonds using the model.

\section{Affine Multi-factor models}
We first construct the multi-factor models. Let $\boldsymbol{x}_{t}=\left(x_{1 t}, x_{2 t}, \ldots,
x_{n t}\right)^{T}$ be a vector of state variables, and assume that these have affine $\Q$-dynamics
given by: 
\[
  \d  \boldsymbol{x}_{t}=
  \left(\underset{n\times 1}{\hat{\varphi}}-\underset{n\times n}{\hat{\kappa}} \underset{n \times
  1}{\boldsymbol{x}_{t}}\right) \d  t+\underset{n \times  n}{\Gamma }\sqrt{\underset{n\times
n}{\boldsymbol{V}\left(\boldsymbol{x}_{t}\right)}} \d  \underset{n \times
1}{\boldsymbol{z}_{t}^{\Q}}
,\]
where $\boldsymbol{V}\left(\boldsymbol{x}_{t}\right)$ is the diagonal matrix
\[
\boldsymbol{V}\left(\boldsymbol{x}_{t}\right)=\left[\begin{array}{cccc}
v_{0,1}+\mathbf{v}_{1,1}^{T} \boldsymbol{x}_{t} & 0 & \cdots & 0 \\
0 & v_{0,2}+\mathbf{v}_{1,2}^{T} \boldsymbol{x}_{t} & \cdots & 0 \\
\vdots & \vdots & \ddots & 0 \\
0 & 0 & \cdots & v_{0, n}+\mathbf{v}_{1, n}^{T} \boldsymbol{x}_{t}
\end{array}\right]
\]
For admissibility we require $v_{0, i}+\mathbf{v}_{1, i}^{T} \boldsymbol{x}_{t} \geq 0$ for all $i$
to ensure that the square root is well-defined, and $\boldsymbol{x}_{t} \in \mathbb{R}^{n}$. We
further assume that the short rate is affine in the state variables, $\bm{x}_t$:
\[
r_{t}=\xi_{0}+\xi^{T} \boldsymbol{x}_{t}
.\]
And we assume that the market price of risk is given by
\[
\lambda\left(\boldsymbol{x}_{t}\right)=\sqrt{\boldsymbol{V}\left(\boldsymbol{x}_{t}\right)} \bar{\lambda}
,\]
for some vector $\overline{\lambda} \in \R^{n}$. 
Such that the $\P$-dynamics take the (same) affine form:
\[
\d  \boldsymbol{x}_{t}=\left(\varphi-\kappa \boldsymbol{x}_{t}\right) \d  t+\Gamma \sqrt{\boldsymbol{V}\left(\boldsymbol{x}_{t}\right)} \d  \boldsymbol{z}_{t}
\]
where $\varphi=\hat{\varphi}+\Gamma \psi, \kappa=\hat{\kappa}-\Gamma \Psi$ with
\[
\psi=\left[\begin{array}{c}
\bar{\lambda}_{1} v_{0,1} \\
\bar{\lambda}_{2} v_{0,2} \\
\vdots \\
\bar{\lambda}_{n} v_{0, n}
\end{array}\right], \quad \text { and } \quad \Psi=\left[\begin{array}{c}
\bar{\lambda}_{1} \mathbf{v}_{1,1}^{T} \\
\bar{\lambda}_{2} \mathbf{v}_{1,2}^{T} \\
\vdots \\
\bar{\lambda}_{n} \mathbf{v}_{1, n}^{T}
\end{array}\right]
\]

\section{Pricing and characterization}
In the model setup above, the price of a zero-coupon bond is given by
\[
B_{t}^{T}=B^{T}\left(x_{t}, t\right)=\e^{-a(T-t)-b(T-t)^{T} x_{t}}
\]
where the functions $a(\tau)$ and $b(\tau)$ solve the Riccati ODE's, i.e. $a(0)=b(0)=0$ and
\[
\begin{aligned}
& b^{\prime}(\tau)=-\hat{\kappa}^{T} b(\tau)-\frac{1}{2} \sum_{i=1}^{n}\left[\Gamma^{T} b(\tau)\right]_{i}^{2} v_{1, i}+\xi \\
& a^{\prime}(\tau)=\hat{\varphi}^{T} b(\tau)-\frac{1}{2} \sum_{i=1}^{n}\left[\Gamma^{T} b(\tau)\right]_{i}^{2} v_{0, i}+\xi_{0}
\end{aligned}
\]
\begin{note*}
There are $n+1$ ODE's, to solve these first find $b(\tau)$ from the first equality and then
integrate the second to obtain $a(\tau)$.
\end{note*}


Affine Multi-factor models of size $n \in \N$ can be divided into $n+1$ categories, denoted by
$\mathbb{A}_m(n)$ for each $m \in\{0,1, \ldots, n\}$. For a fixed $m$ each class is characterized
by:
\begin{itemize}
  \item The first $m$ factors affect the volatility $\boldsymbol{V}\left(\boldsymbol{x}_{t}\right)$, and are known as volatility factors.
  \item The remaining $n-m$ do not affect $\boldsymbol{V}\left(\boldsymbol{x}_{t}\right)$ and are known as non-volatility factors. Vasicek is $\mathbb{A}_{0}(1)$ model, CIR is $\mathbb{A}_{1}(1)$. We need to be careful about the admissibility requirement.
\end{itemize}
\subsubsection{Parameter restrictions for volatility factors}
For a volatility factor $x_i$ where $i \in 1,\ldots ,m$, to ensure admissibility we need $x_i$ to be
non-negative, hence we impose the following conditions:
\begin{itemize}
  \item $\varphi_{i}>0, \quad \kappa_{i i}>0, \quad$ \text{for mean reversion,}
  \item $\kappa_{i j}=0$, for $j=m+1, \ldots, n, \quad \kappa_{i j} \leq 0$, for $j=1, \ldots, m$,
  with $j \neq i$,
\item Furthermore, $x_{i}$ must only be sensitive to $\text{d} z_{i t}$ with volatility coefficient proportional to $\sqrt{x_{i t}}$.
\end{itemize}
This leads to
\[
\d  x_{i t}=\left(\varphi_{i}-\sum_{j=1}^{n} \kappa_{i j} x_{j t}\right) \d  t+\Gamma_{i i}
\sqrt{x_{i t}} \d  z_{i t} .
\]
Thus, volatility factors are instantaneously uncorrelated square root processes. But they are
not necessarily independent because volatility factors affect each others drifts (positively). Also,
the volatility factors cannot be dependent on the non-volatlity factors (they may be negative).

\subsubsection{Parameter restrictions for non-volatility factors}
For the non-volatility factors $x_i$, to ensure admissibility:
\begin{itemize}
  \item Drift rate may depend on all factors.
  \item The sensitivity to the volatility shocks $\d  z_{j t}$ must be proportional to $\sqrt{x_{j
    t}}\; (j= 1,2, \ldots, m)$.
  \item The sensitivity to non-volatility shocks $\d  z_{j t}\; (j=m+1, \ldots, n)$ is
    \[
      \Gamma_{i i} \sqrt{v_{0, j}+\sum_{k=1}^{m} v_{1, j k} x_{k t}} \text {, }
    \]
    with $v_{0, j}, v_{1, j k} \geq 0$.
\end{itemize}
This leads to
\[
\d  x_{i t}=\left(\varphi_{i}-\sum_{j=1}^{n} \kappa_{i j} x_{j t}\right) \d  t+\sum_{j=1}^{m}
\Gamma_{i j} \sqrt{x_{j t}} \d  z_{j t}+\sum_{j=m+1}^{n} \Gamma_{i j} \sqrt{v_{0, j}+\sum_{k=1}^{m}
v_{1, j k} x_{k t}} \d  z_{j t} .
\]

\subsection{Relevant results}

Principal component analysis with US bond data shows
that first factor (level) explains about $85 \%$ of the variation in yields, second factor (slope)
about $10 \%$, third factor (curvature) about $2 \%$. So we need at least two factors. So generally,
there is marginal return in model accuracy compared to complexity if we add more than three
factors.

\begin{description}
\item[Generalized affine models] 
\end{description}
 The idea is to allow a general market price of risk, which gives more flexible setting when
 estimating model parameters

 \begin{description}
   \item[Completely affine models]
 \end{description}
\[
\begin{aligned}
r_{t} & =\xi_{0}+\xi^{T} x_{t} \\
\d  x_{t} & =\left(\hat{\varphi}-\hat{\kappa} x_{t}\right) \d  t+\Gamma \sqrt{\boldsymbol{V}\left(x_{t}\right)} \d  z_{t}^{\Q} \\
\lambda\left(x_{t}\right) & =\sqrt{\boldsymbol{V}\left(x_{t}\right)} \bar{\lambda}
\end{aligned}
\]
The market price of risk is affected only by volatility factors.

\begin{description}
  \item[ Essentially affine models]
\end{description}
\[
\lambda_{t}=\sqrt{\boldsymbol{V}\left(x_{t}\right)} \bar{\lambda}_{1}+\sqrt{\boldsymbol{V}\left(x_{t}\right)^{-}} \bar{\lambda}_{2} x_{t}
\]
where $\boldsymbol{V}\left(x_{t}\right)^{-}$is diagonal matrix (broadly speaking the inverse of
$\boldsymbol{V}\left(x_{t}\right)$ ). Here, the non-valotility factors can also affect the market
price of risk.

\begin{description}
  \item[Extended affine models]
\end{description}
\[
\sqrt{\boldsymbol{V}\left(x_{t}\right)} \lambda_{t}=\bar{\lambda}_{1}+\bar{\lambda}_{2} x_{t}
\]
with restrictions in $\bar{\lambda}_{2}$ to ensure that the process for $x_{t}$ is well-defined
under $\Q$ and $\P$. Offers more flexibility for the volatility factors on the market price of
risk.

